\documentclass[11pt,a4paper,oneside,openany,parskip=false,bibliography=totoc,numbers=noenddot]{scrbook}
% ---------------------------------------------------------------- page
\usepackage[T1]{fontenc}
\usepackage[utf8]{inputenc}
\usepackage[bahasai]{babel}
\usepackage{amsmath,mathtools,bm}
\usepackage{newpxtext,newpxmath}
\usepackage[scaled=0.93]{sourcesans}
\usepackage[a4paper,top=28mm,bottom=30mm,inner=32mm,outer=32mm,headsep=10mm]{geometry}
\usepackage{microtype}
\linespread{1.15}
\setlength{\parindent}{1.2em}
\setlength{\parskip}{0pt}

% ---------------------------------------------------------------- tables, figures, lists
\usepackage{booktabs,tabularx,array,multirow}
\usepackage{graphicx}
\usepackage{enumitem}
\setlist{itemsep=1pt,topsep=4pt,leftmargin=1.6em}
\usepackage{xcolor}
\usepackage{tikz,pgfplots}
% Shared TikZ setup. Used by the book and by tools/tikz2svg.py (standalone figures for the web).
\usetikzlibrary{arrows.meta,positioning,calc,shapes.geometric,shapes.misc,decorations.pathreplacing,fit,backgrounds,patterns,matrix}
\usepgfplotslibrary{fillbetween}
\pgfplotsset{compat=1.18}
\definecolor{biru}{HTML}{1F5A8C}
\definecolor{hijau}{HTML}{1E7A5A}
\definecolor{oranye}{HTML}{C0661E}
\definecolor{ungu}{HTML}{6A3D9A}
\definecolor{merah}{HTML}{B03A3A}
\definecolor{abu}{HTML}{6B6B6B}
\definecolor{pasir}{HTML}{F4EDE2}
\tikzset{
  >={Stealth[length=2.2mm]},
  kotak/.style={draw=#1!80!black,fill=#1!8,rounded corners=2pt,minimum height=7mm,minimum width=14mm,align=center,font=\sffamily\small},
  kotak/.default=biru,
  bulat/.style={draw=#1!80!black,fill=#1!10,circle,minimum size=7mm,inner sep=1pt,font=\sffamily\small},
  bulat/.default=biru,
  panah/.style={->,thick,draw=abu},
  label/.style={font=\sffamily\footnotesize,text=abu},
}

\usepackage[font=small,labelfont=sc,labelsep=period,margin=8mm]{caption}

% ---------------------------------------------------------------- citations, links
\usepackage[round,authoryear,sort]{natbib}
\setlength{\bibsep}{3pt}
\usepackage[noautomatic]{imakeidx}
\makeindex[title=Indeks,columns=2,intoc]
\usepackage[hidelinks,unicode]{hyperref}
\usepackage{xurl}
\usepackage[nameinlink,noabbrev]{cleveref}
\crefname{chapter}{Bab}{Bab}
\crefname{section}{Subbab}{Subbab}
\crefname{figure}{Gambar}{Gambar}
\crefname{table}{Tabel}{Tabel}
\crefname{equation}{Persamaan}{Persamaan}
\crefname{part}{Bagian}{Bagian}
\crefname{appendix}{Lampiran}{Lampiran}
\Crefname{appendix}{Lampiran}{Lampiran}
\newcommand{\crefpairconjunction}{ dan }
\newcommand{\crefmiddleconjunction}{, }
\newcommand{\creflastconjunction}{, dan }
\newcommand{\crefrangeconjunction}{ sampai }
\creflabelformat{equation}{#2(#1)#3}

% ---------------------------------------------------------------- headings (KOMA), classic serif
\setkomafont{disposition}{\normalfont\rmfamily}
\addtokomafont{chapter}{\huge}
\addtokomafont{part}{\Huge}
\addtokomafont{section}{\Large}
\addtokomafont{subsection}{\large\itshape}
\addtokomafont{subsubsection}{\normalsize\itshape}
\renewcommand*{\chapterformat}{\thechapter\enskip}
\renewcommand*{\chapterheadstartvskip}{\vspace*{14mm}}
\renewcommand*{\chapterheadendvskip}{\vspace{10mm}}
\RedeclareSectionCommand[beforeskip=-3ex plus -1ex,afterskip=1.2ex]{section}
\RedeclareSectionCommand[beforeskip=-2.2ex plus -0.6ex,afterskip=0.8ex]{subsection}
\usepackage[automark]{scrlayer-scrpage}
\clearpairofpagestyles
\ohead{\small\scshape\MakeLowercase{\leftmark}}
\ofoot*{\small\pagemark}
\setkomafont{pageheadfoot}{\normalfont}
\renewcommand*{\chaptermarkformat}{}

% A short note on where a chapter's material comes from, set at the end of the chapter.
\newcommand{\sumberbab}[1]{\par\bigskip\noindent{\small\textsc{Catatan sumber.}\enskip #1\par}}

% ---------------------------------------------------------------- notation
\newcommand{\R}{\mathbb{R}}
\newcommand{\E}{\mathbb{E}}
\newcommand{\Prob}{\mathbb{P}}
\newcommand{\N}{\mathcal{N}}
\newcommand{\Loss}{\mathcal{L}}
\newcommand{\dd}{\,\mathrm{d}}
\newcommand{\vx}{\bm{x}}
\newcommand{\vy}{\bm{y}}
\newcommand{\vz}{\bm{z}}
\newcommand{\vw}{\bm{w}}
\newcommand{\vu}{\bm{u}}
\newcommand{\vh}{\bm{h}}
\newcommand{\vtheta}{\bm{\theta}}
\newcommand{\mW}{\bm{W}}
\newcommand{\mA}{\bm{A}}
\newcommand{\mX}{\bm{X}}
\newcommand{\KL}{\mathrm{KL}}
\newcommand{\ket}[1]{\lvert #1\rangle}
\newcommand{\braket}[2]{\langle #1\vert #2\rangle}
\DeclareMathOperator*{\argmin}{arg\,min}
\DeclareMathOperator*{\argmax}{arg\,max}
\newcommand{\istilah}[1]{\emph{#1}\index{#1}}


\newcommand{\bab}[1]{\IfFileExists{chapters/#1.tex}{\input{chapters/#1}}{}}

\begin{document}
\frontmatter
\begin{titlepage}
\thispagestyle{empty}
\begin{tikzpicture}[remember picture,overlay]
  \fill[pasir] (current page.south west) rectangle (current page.north east);
  % thin rule that anchors the text block on the left
  \draw[biru!80!black,line width=2.2pt]
    ($(current page.north west)+(2.2,-3.2)$) -- ($(current page.north west)+(2.2,-11.3)$);
  % a quiet motif: scattered points that slowly organise into a flow
  \begin{scope}[shift={($(current page.south west)+(2.2,2.4)$)}]
    \foreach \i in {0,...,16}{
      \foreach \j in {0,...,6}{
        \pgfmathsetmacro{\yy}{\j*0.85+0.35*sin(\i*25+\j*40)*(1-\i/18)}
        \pgfmathsetmacro{\op}{0.12+0.045*\i}
        \fill[biru,opacity=\op] (\i*1.0,\yy) circle (0.05+0.004*\i);
      }
    }
    \draw[oranye,line width=1.1pt,opacity=0.85]
      plot[smooth,domain=0:16,samples=80] (\x,{2.6+1.1*sin(\x*28)*exp(-0.05*\x)});
  \end{scope}
  % text block, aligned on one left edge
  \node[anchor=north west,align=left,text width=13.5cm,inner sep=0pt]
    at ($(current page.north west)+(2.75,-3.2)$) {%
    {\large\scshape\color{abu} catatan kuliah master artificial intelligence di linz\par}
    \vspace{9mm}
    {\fontsize{46}{52}\selectfont\color{biru!80!black} Belajar\\[1mm] Mesin\hspace{0.3em}Belajar\par}
    \vspace{8mm}
    {\Large\itshape\color{oranye!80!black} dari data dan neuron sampai sinyal,\\ molekul, dan fluida\par}
    \vspace{16mm}
    {\large Eigenzayn\par}
    \vspace{1.5mm}
    {\small\color{abu} Edisi pertama, Oktober 2026\par}};
\end{tikzpicture}
\end{titlepage}

\thispagestyle{empty}
\vspace*{\fill}
\noindent{\small\color{abu}
\textit{Belajar Mesin Belajar}\\
Catatan kuliah Master Artificial Intelligence di Linz\\[3mm]
Isi buku ini ditulis ulang dengan kata-kata sendiri dari kuliah yang diikuti, lalu diperluas
dengan bacaan di luar kuliah: bagian perkembangan riset di akhir bab, beberapa bab yang disusun
dari literatur, dan lampiran. Slide, naskah, dan soal ujian asli tetap milik para pengajarnya
dan tidak disertakan. Semua rujukan di daftar pustaka diperiksa terhadap Crossref, arXiv, atau
halaman hak cipta bukunya. Laporan, tanggal, dan angka peristiwa terbaru diperiksa dengan
sumber resminya pada saat penulisan.\\[3mm]
Disusun dengan \LaTeX.
}
\clearpage

\chapter*{Prakata}
\addcontentsline{toc}{chapter}{Prakata}

Pada semester musim dingin 2024/25 saya mulai belajar kecerdasan buatan di Linz. Semester
pertama langsung padat: computer vision, jaringan saraf, dan LSTM yang diajarkan oleh
penemunya sendiri. Setiap kuliah meninggalkan setumpuk slide. Dua tahun kemudian, map di
laptop berisi ratusan PDF.

Slide bagus untuk mengikuti kuliah, tetapi kurang enak dibaca ulang. Kalimatnya terpotong,
rumus muncul tanpa cerita, dan hubungan antar mata kuliah jarang terlihat. Padahal hubungan
itulah yang paling menarik. Konvolusi di kuliah computer vision ternyata saudara dekat
stensil beda hingga di kuliah mekanika fluida. Model difusi di kuliah machine learning lanjut
ternyata bersandar pada persamaan Fokker--Planck dari fisika statistik. Kalman filter di kuliah
model probabilistik ternyata cara yang sama yang dipakai ahli cuaca untuk menggabungkan
simulasi dengan pengukuran.

Buku ini adalah catatan perjalanan itu. Saya menulis ulang isi kuliah dengan bahasa sendiri,
menambahkan gambaran yang dulu membantu saya memahaminya, lalu merujuk literatur yang lebih
luas dari sekadar slide. Harapannya, siapa pun yang membaca, entah mahasiswa, insinyur, atau
orang yang sekadar penasaran, bisa mengikuti ceritanya dari satu neuron sampai model yang
meniru aliran fluida.

Di tengah perjalanan, satu pertanyaan semakin sering muncul di kepala saya: bagaimana machine
learning bisa membantu simulasi fisika, dan bagaimana fisika bisa membuat machine learning lebih
bisa dipercaya? Pertanyaan itu menjadi benang merah buku ini. Hampir setiap bab ditutup dengan
melihat materinya dari sudut simulasi, dan Bagian III membahas dunia \emph{physics-informed
machine learning} secara utuh: PINN, neural operator, simulator berbasis graf, penemuan
persamaan dari data, dan model hibrida.

Gaya penulisannya sengaja dibuat seperti catatan yang dibaca, bukan buku latihan. Setiap
gagasan besar diberi satu gambaran dari kehidupan sehari-hari sebelum rumusnya muncul. Analogi
tidak pernah sempurna, dan kalau sebuah analogi mulai menyesatkan, teks akan mengatakannya.
Contoh hitungan kecil ditulis di dalam paragraf, supaya pembaca bisa mengikutinya dengan
pensil kalau mau. Istilah teknis dibiarkan dalam bahasa Inggris dan dicetak miring saat pertama
muncul, karena istilah itulah yang dipakai di kuliah dan di makalah.

Setiap bab diakhiri catatan singkat tentang mata kuliah asal dan bacaan utamanya. Slide dan
soal asli tetap milik para pengajarnya dan tidak disertakan. Yang ada di sini adalah pemahaman
saya atas materi itu, jadi kekeliruan yang tersisa adalah kekeliruan saya.

\bigskip
\hfill Linz, Oktober 2026

\setcounter{tocdepth}{0}
\tableofcontents

\mainmatter
\setcounter{chapter}{-1}
\bab{ch00_peta}

\part{Fondasi}
\bab{ch01_belajar_dari_data}
\bab{ch01b_statistik_bayes}
\bab{ch01c_teori_belajar}
\bab{ch02_optimasi}
\bab{ch02b_ai_klasik}
\bab{ch02c_unsupervised}

\part{Deep Learning dan Persepsi}
\bab{ch03_dl_basic}
\bab{ch03b_praktik}
\bab{ch04_lstm_transformer}
\bab{ch04b_nlp}
\bab{ch05_computer_vision}
\bab{ch06_arsitektur_generatif}
\bab{ch07_geometric}
\bab{ch08_probabilistik}
\bab{ch09_ml_advanced}
\bab{ch10_rl}
\bab{ch11_xai}

\part{Sinyal, Sensor, dan Sistem}
\bab{ch30_sinyal}
\bab{ch31_kendali}
\bab{ch32_pervasif}
\bab{ch33_robot}

\part{AI untuk Ilmu Hayati}
\bab{ch40_genomik}
\bab{ch41_molekul}
\bab{ch42_protein}
\bab{ch43_medis}

\part{AI untuk Simulasi}
\bab{ch12_persamaan_fluida}
\bab{ch13_numerik_fluida}
\bab{ch14_fisika_komputasi}
\bab{ch15_pinn}
\bab{ch16_operator}
\bab{ch17_penemuan_reduksi}
\bab{ch18_hibrida}
\bab{ch19_rcfd}

\part{AI, Manusia, dan Masa Depan}
\bab{ch50_rekomendasi}
\bab{ch51_ai_manusia}
\bab{ch51c_keselamatan}
\bab{ch51b_kuantum}
\bab{ch52_terbaru}

\appendix
\bab{lampA_glosarium}
\bab{lampB_matematika}
\bab{lampC_riset}

\backmatter
\bibliographystyle{plainnat}
\bibliography{refs}
\printindex
\end{document}
